\documentclass[11pt]{article}

\usepackage[margin=25mm]{geometry}

\usepackage{amsmath,amssymb,fontspec,unicode-math}

\IfFontExistsTF{DejaVu Serif}{\setmainfont{DejaVu Serif}}{\setmainfont{Latin Modern Roman}}

\IfFontExistsTF{Latin Modern Math}{\setmathfont{Latin Modern Math}}{\setmathfont{DejaVu Math TeX Gyre}}

\usepackage{xurl,hyperref}

\setlength{\parindent}{0pt}\setlength{\parskip}{6pt}

\begin{document}

\title{HETEROGENEOUS SWARMS: Jointly Optimizing Model Roles and Weights for Multi-LLM Systems}

\author{ELVIA preliminary AI analysis}\date{21 October 2026}\maketitle

Preliminary assessment: Minor.\par

Independent scientific review remains pending. Author adoption has not been requested.

The latest author port corrects the earlier final-answer evaluator and unstable softmax. Two bounded concerns remain: graph decoding mutates an adjacency row and restricts the graph family, while the current paper’s finite-sampling contribution score is undefined when a model is never sampled. Neither establishes damage to the reported aggregate benchmark results.

\section*{Review scope}

Latest-version, claim-focused review of the cited author-paper passages and applicable current author-release paths. The stated counterexamples and scope controls were checked analytically; no author program, old verifier, training run or benchmark was executed.

\section*{Problem 1: Decoder topology and input mutation differ from the intended reusable DAG representation}

Source: Section 3 graph representation; current MoCo model\_collaboration/method/text\_heterogeneous\_swarms.py, graph\_decode

Assessment: supported; impact: minor

\subsection*{Source claim (paraphrase)}

The paper presents general DAG collaboration; the release attaches each newly selected node to one existing node and reuses adjacency rows.

\subsection*{Analysis}

A successful \(n\)-node decode adds \(n-1\) edges with one successor per nonterminal node. It cannot represent every DAG, for example \(0\to1,0\to2,1\to2\). The current author port still assigns adjacency\_matrix[selected\_node] to a working row without copying it. NumPy row indexing returns a view, so zeroing entries for the unselected nodes changes the caller’s matrix. This is a restricted graph family and a side effect; the output remains acyclic, and corruption of every optimizer state is not established.

\subsection*{Finding and scope}

The same two local properties remain in the latest author-directed MoCo port. Their effect on saved optimization trajectories or published scores has not been established.

\subsection*{Proposed correction}

Copy the row before mutation. That repairs the side effect but does not enlarge the graph family; intended edge-sampling semantics remain to be specified.

The author’s original repository directs readers to MoCo. The current MoCo dev branch retains the row view and single-successor construction; no correction of those properties was found in the inspected function.

\section*{Problem 2: The latest evaluator corrects the historical numeric substring mismatch}

Source: Current MoCo model\_collaboration/data/eval.py, extract\_answer\_text, calculate\_exact\_match and get\_scores

Assessment: not\_issue; impact: no\_issue\_found

\subsection*{Source claim (paraphrase)}

The original evaluator treated inclusion of the reference string as an exact match and had a separate judge branch with different score units. The latest author port extracts the final answer and uses type-appropriate equality or numeric comparison.

\subsection*{Analysis}

For the GSM8K example, the correct answer is \(14(6+2)=112\), whereas the wrong answer \(1120\) contains its decimal substring. That distinguishes the old substring rule from exact final-answer comparison. In the current calculate\_exact\_match, normalized numeric strings are parsed and checked by \_numbers\_match; these two numbers fail its relative tolerance of \(10^{-2}\). Text falls back to normalized equality after final-answer extraction. The earlier integer-judge branch is not used by the current swarm method’s get\_scores path.

\subsection*{Finding and scope}

The precise historical false-positive witness is corrected in the latest author implementation and is not counted as a current problem. This does not establish that historical predictions or tables were recomputed.

\section*{Problem 3: Finite random assignments can leave JFK-score undefined}

Source: Current paper Algorithm 2, PDF p. 5; current MoCo text\_heterogeneous\_swarms.py, run\_method

Assessment: supported; impact: minor

\subsection*{Source claim (paraphrase)}

Algorithm 2 divides a model’s summed sampled utility by the number of appearances in a finite set of IID assignments, without requiring positive coverage of every model.

\subsection*{Analysis}

With \(10\) models, \(10\) assignments and \(10\) positions per assignment, there are \(100\) IID model draws. A specified model is absent with probability \(0.9^{100}>0\). For example, appearance counts \[(12,16,15,0,7,9,9,9,8,15)\] sum to \(100\) but give the absent model utility and frequency both zero. Its prescribed ratio is therefore \(0/0\), even for constant positive assignment utility. The latest MoCo port instead fixes model\_names and optimizes adjacency; it omits this frequency-based role-selection stage. That removal prevents attributing the same division to the current port, but does not define the score in the still-current published Algorithm 2.

\subsection*{Finding and scope}

The current paper’s finite-sampling definition lacks a zero-frequency convention. This is a domain defect with a positive-probability event, not evidence that a reported run encountered it or that the current MoCo port performs the same division.

\subsection*{Proposed correction}

Define behavior for zero-frequency models. Conditioning samples on complete coverage makes denominators positive but changes the sampling law; no unbiasedness equivalence is asserted.

The current paper retains the unguarded finite-sampling score. The later MoCo port omits that role-selection stage rather than providing a corrected implementation of its stated ratio.

\section*{Problem 4: The latest stable softmax corrects the historical overflow witness}

Source: Current MoCo text\_heterogeneous\_swarms.py, softmax and graph\_decode

Assessment: not\_issue; impact: no\_issue\_found

\subsection*{Source claim (paraphrase)}

The original decoder exponentiated reciprocal row sums directly. The latest author port subtracts their finite maximum and separately handles masked or nonfinite entries.

\subsection*{Analysis}

Positive row weights can yield reciprocal row sums \((5000,1,1)\). Direct float64 exponentiation of \(e^{5000}\) overflows despite a well-defined real softmax. The current implementation calculates exponentials after subtracting the maximum finite value, so this witness becomes \((e^0,e^{-4999},e^{-4999})\) before normalization. Its code also handles positive infinity and masks explicitly. This is the previously identified issue and later author correction, not a novelty claim.

\subsection*{Finding and scope}

The exact historical overflow defect is corrected in the latest author port. No general claim is made about every possible nonfinite input or downstream numerical operation.

\section*{Problem 5: Table 7 supports a mean advantage but not every-task superiority}

Source: Appendix D, Table 7, PDF p. 28.

Assessment: not\_issue; impact: no\_issue\_found

\subsection*{Source claim (paraphrase)}

The text says the mixed \(7\)-billion/\(2\)-billion-parameter system exceeds the best single \(7\)-billion-parameter model, without explicitly claiming superiority on every task.

\subsection*{Analysis}

The four mixed-minus-single scores are \[(-0.028,0.134,-0.043,0.029)\], with arithmetic means \(0.31875\) and \(0.29575\). Thus the aggregate advantage is consistent with the table, while two task reversals remain. The rounded aggregate arithmetic in Table 1 is also consistent.

\subsection*{Finding and scope}

This is a scope clarification, not an additional demonstrated error. Table 1’s rounded \(18.5\%\) aggregate remains a positive arithmetic control.

\section*{Limits}

No author programs, optimizers or model evaluations were executed. Historical benchmark predictions and run commits were unavailable. The current MoCo port changes the role-selection procedure, so absence of its frequency division is not a validation of the published Algorithm 2.\par

\section*{Source identity}

NeurIPS 2025 official proceedings PDF (current publisher bytes, checked 2026-10-21)

\url{https://proceedings.neurips.cc/paper_files/paper/2025/file/a645fda7a838e3cf96ee2bfd0421ca74-Paper-Conference.pdf}

Reviewed PDF SHA-256: \nolinkurl{f44c084ff7c9fdf18b8d2275a1dc68d1af16934bbbb0eff034427bd36e0e7c03}

This is the same research work as the predecessor, matched by title, author identity and official publisher channel. Its freshly retrieved exact PDF bytes and real retrieval time are recorded. Current target arXiv metadata and the relevant available author-version passages were checked separately; background works were excluded. Original report/source history is unchanged.

\end{document}
